\chapter*{如何使用本书}
\addcontentsline{toc}{chapter}{如何使用本书}

本书既可以作为一门研究生课程的教材，也可以作为开展截断维格纳近似（以下简称 TWA）数值研究时的案头参考。全书的主线不是罗列相空间公式，而是反复回答同一组问题：量子态如何变成可抽样的初始分布，算符演化如何变成相空间轨迹，轨迹平均如何还原实验观测量，以及近似和数值离散在什么条件下仍然可信。读者在进入某一章以前，最好先明确自己要解决的是上述哪一个环节。

\section*{读者需要具备什么基础}

阅读本书需要熟悉本科量子力学中的态矢、密度算符、对易关系和谐振子，并能够使用傅里叶变换、常微分方程和基本概率统计。了解二次量子化、玻色凝聚和 Gross--Pitaevskii 方程会使第三、四部分更容易，但第10章仍从有限模式展开重新建立所需记号。涉及路径积分、随机微分方程或 Bogoliubov 理论的内容均在首次使用时给出必要背景。

读者不必在开始阅读前熟练掌握编程。不过，若希望复现数值实验，应具备基本的 Python 数组运算和绘图经验，并理解步长、网格、样本数与随机种子各自控制什么误差。附录C给出数值算法和测试规范，可以在阅读第12章前提前查阅。

\section*{五种建议阅读路线}

各章之间存在明确的依赖关系，但并非每位读者都需要从头顺序读到结尾。可以根据目标选择下列路线。

\begin{description}[style=nextline,leftmargin=2.2em,font=\bfseries]
  \item[系统学习路线]
  按照第1--18章顺序阅读。第一部分建立 Weyl--Wigner 表示，第二部分说明 TWA 从何而来以及为何失效，第三部分把单模方法推广到离散玻色场，第四部分完成双分量淬火案例，第五部分讨论开放系统、替代方法与研究工作流。这条路线适合一学期课程和首次系统接触相空间方法的读者。

  \item[快速开始 TWA 计算]
  先读第1、4、7和9章，掌握问题结构、排序修正、截断步骤和有效性边界；随后完整阅读第10--13章，依次完成模式截断、初态抽样、经典场传播和观测量重构。遇到推导细节时回查第2、3章及附录A。快速路线不能跳过第9章，因为“方程可以运行”并不等于 TWA 在目标时间尺度上有效。

  \item[偏重形式理论与推导]
  重点阅读第2、3、5--9章，并配合附录A、B。建议亲手推导谐振子、Kerr 振子和 Bose--Hubbard 二聚体的 Weyl 符号，再比较精确 Moyal 演化与截断后的 Liouville 流。第5章的二次体系是符号、归一化和传播器实现的基准，第6、8章则用精确小系统暴露非线性近似的误差来源。

  \item[双分量凝聚体与淬火动力学]
  在第1、4、7和9章之后，阅读第10--17章。第14章区分均匀组分耦合与具有动量转移的拉曼自旋轨道耦合；第15、16章研究线性化淬火、相位扩散和模式耦合；第17章集中讨论条纹序、有限尺寸标度及超固态判据。这里必须把动力学中出现条纹、存在相干峰和证明平衡态自发对称性破缺视为不同层次的命题。

  \item[数值实现与研究复现]
  先读第4、7、10--13章和附录B、C，再运行随书程序。对每个案例都应先复现默认基准，然后分别改变样本数、时间步长、空间网格、模式截止和随机种子；只有主要结论在相应收敛区间内稳定，才能把图像解释为物理结果。准备独立研究时，再用第9章和第18章选择精确对照、替代方法及失效诊断。
\end{description}

\section*{每一章应当怎样读}

每章开头的“本章目标”说明读完后应能完成的任务，而不是仅需记住的术语。建议按照下面四步使用一章内容。

\begin{enumerate}[label=\textbf{第\arabic*步},leftmargin=4.2em]
  \item \textbf{先定位问题。} 在阅读推导前写下本章的输入、输出和近似。例如，输入是密度算符还是 Bogoliubov 协方差，输出是单条轨迹还是正规排序关联函数。若无法说清这两点，复杂公式通常也难以正确使用。

  \item \textbf{沿计算链阅读。} 对每个模型追踪
  \[
    \begin{aligned}
      \text{量子态与哈密顿量}
      &\longrightarrow \text{Weyl 符号与初态抽样}\\
      &\longrightarrow \text{轨迹方程}\\
      &\longrightarrow \text{观测量估计器}\\
      &\longrightarrow \text{误差与有效性判断}.
    \end{aligned}
  \]
  不要只抄写中间的运动方程。漏掉初态的半量子噪声或末端的排序修正，都可能得到看似平滑但物理含义错误的结果。

  \item \textbf{完成最小验证。} 对解析章节检查一个可解极限；对数值章节至少检查维度、守恒量、步长收敛、网格或模式截止收敛以及统计误差。书中“数值实验”给出可复现基准，“有效性警告”指出即使程序无报错仍可能出现的系统偏差。

  \item \textbf{用章末问题闭环。} “章末检查”适合口头自测，习题用于补齐推导和实现细节。附录D提供提示与部分答案；建议先保留完整推导和失败尝试，再与提示比较，而不是从答案反推步骤。
\end{enumerate}

若把本书用于课程教学，可以把一次课分成“物理问题与表示”“核心推导”“数值或解析验证”三段。若用于自学，每章结束时应能用一页纸回答三个问题：本章作了什么近似；什么量可以与实验或精确结果比较；哪一个检验最可能推翻当前结论。

\section*{必须始终区分的平均与排序}

本书使用相似的尖括号或平均符号描述不同操作。它们不能互换：
\begin{equation*}
  \text{量子期望}\quad
  \qavg{\Ohat}=\Tr(\rhohat\Ohat),
  \qquad
  \text{Wigner 样本平均}\quad
  \ensavg{F}=\frac{1}{N_{\mathrm{s}}}
  \sum_{s=1}^{N_{\mathrm{s}}}F[\psi_s].
\end{equation*}
空间平均、时间平均和无序平均会在使用它们的章节中另行定义。一条 Wigner 轨迹不是一次完整的量子测量，也通常不等于量子期望；它是用于构造系综统计的相空间样本。

Wigner 平均天然给出对称排序量，而粒子数、密度和计数关联通常采用正规排序。最简单的单模例子是
\begin{equation*}
  \qavg{\had\ha}=\ensavg{|\alpha|^2}-\frac12.
\end{equation*}
在有限模式场中，$1/2$ 被局域投影核 $\delta_{\mathcal C}(x,x)/2$ 取代。因而，读者在记录任何数值观测量时，都应同时记录它的算符定义、排序约定和相空间估计器。第4章介绍单模修正，第10、13章给出场论版本，附录A汇总常用公式。

\section*{如何使用随书数值材料}

获得本书源码后，应从项目根目录运行程序。各章程序位于 \path{code/}，轻量数值结果位于 \path{data/}，生成的图位于 \path{figures/}。命令
\begin{lstlisting}[style=bookpython,numbers=none]
python code/run_all.py
\end{lstlisting}
会依次运行书中的九个数值基准，并把通过状态、运行时间和输出摘要写入 \path{data/test_summary.json}。单个脚本也可以独立运行，适合在修改参数后快速检查。实际 Python 可执行文件和依赖版本应由读者自己的环境记录；不能仅凭一张相似图像认定复现成功。

建议为每项研究建立四层记录：第一层保存模型参数、单位和随机种子；第二层保存代码版本及运行环境；第三层保存原始轨迹或足以重新计算观测量的数据；第四层保存制图脚本和最终统计量。正文图像用于说明物理结论，不替代原始数据和收敛证据。若修改算法、截止或估计器，应把修改后的结果与默认基准并列，而不是覆盖基准。

\section*{把本书方法用于新问题之前}

从本书案例迁移到新体系时，建议先回答以下问题：保留了哪些模式，舍弃了哪些模式；初态 Wigner 函数能否可靠抽样；Moyal 方程中删去了哪些高阶导数；占据数和演化时间是否处于受控范围；正规排序修正相对于物理信号有多大；结果是否通过样本、步长、网格和截止扫描；是否存在精确小系统、Bogoliubov 理论或其他多体方法可作交叉验证。

这些问题中任何一个没有答案，都不必立刻否定 TWA，但应降低结论强度并明确标注限制。尤其要避免三种常见推断：把单条轨迹的结构当作量子系综序，把有限尺寸中的峰当作热力学极限相变，以及把教学模型的参数趋势直接当作具体实验的定量预言。第9章提供系统的有效性审计，第17章讨论序参量和有限尺寸证据，第18章说明何时应改用正-$P$、张量网络、精确对角化或开放系统方法。

最后，本书所有连续场表达式都应理解为先在有限模式子空间或有限网格上定义，再讨论改变截止时的行为。增加网格点会同时增加真空半量子噪声，并不自动产生更精确的连续极限。读者若始终保留“表示---传播---估计---验证”这条主线，就能把本书当作一套可审计的研究流程，而不只是一组可调用的公式。
